
\SGAStatement{Corollary 3.7}{A unipotent algebraic group \(G\) is
nilpotent (Exposé VI\textsubscript{B}, \S8); its Lie algebra
\(\mathfrak g\) is nilpotent (Bourbaki, \emph{Groupes et algèbres de
Lie}, Chap.~I, \S4) and is isomorphic to a Lie algebra of nilpotent
endomorphisms of a finite-dimensional vector space.  In
characteristic \(p>0\), \(\mathfrak g\) is a unipotent \(p\)-Lie
algebra (\SGARefLink{sga3:xvii:definition:3:6}{3.6}).}
\SGARefTarget{sga3:xvii:corollary:3:7}

Since \SGARefLink{sga3:xvii:theorem:3:5:item:i}{(i)} implies (v), it is enough to prove \SGARefLink{sga3:xvii:corollary:3:7}{3.7} when
\(G=\operatorname{Trigstr}(n)\).  We have already used the fact that
\(\operatorname{Trigstr}(n)\) is a nilpotent algebraic group.
Moreover, the Lie algebra \(\mathfrak h\) of
\(\operatorname{Trigstr}(n)\) consists of the upper-triangular
endomorphisms of \(V\) that have zeros on the main diagonal.  These
endomorphisms are therefore nilpotent, and consequently
\(\mathfrak h\) is nilpotent (Bourbaki, \emph{loc.\ cit.},
Chap.~I, \S4, Cor.~3).  If \(p>0\), the \(p\)-th power in the
\(p\)-Lie algebra
\(\mathfrak{gl}(V)=\operatorname{End}(V)\) coincides with the
\(p\)-th power of endomorphisms of \(V\)
(Exposé VII\textsubscript{A}, \SGARefLink{sga3:VIIA:6.4.4}{6.4.4}); hence \(\mathfrak h\) is
unipotent.

\SGAStatement{Corollary 3.8}{Let \(k\) be an algebraically closed
field, and let \(G\) be a smooth affine algebraic group over \(k\).
Then the following properties are equivalent:
\begin{enumerate}
\item[(i)] \SGARefTarget{sga3:xvii:corollary:3:8:item:i}
  \(G\) is unipotent.
\item[(ii)] \SGARefTarget{sga3:xvii:corollary:3:8:item:ii}
  \(G(k)\) consists of unipotent elements
  (BIBLE~4, Prop.~4, Cor.~1); that is, \(G\) is unipotent in the
  sense of BIBLE.
\end{enumerate}}
\SGARefTarget{sga3:xvii:corollary:3:8}

The implication \SGARefLink{sga3:xvii:corollary:3:8:item:i}{(i)}\(\Rightarrow\)\SGARefLink{sga3:xvii:corollary:3:8:item:ii}{(ii)} follows because, by
\SGARefLink{sga3:xvii:proposition:3:2}{3.2}, \SGARefLink{sga3:xvii:corollary:3:8:item:i}{(i)}\(\Rightarrow\)(v), \(G\) is isomorphic to an algebraic
subgroup of a group \(\operatorname{Trigstr}(n)\).

For \SGARefLink{sga3:xvii:corollary:3:8:item:ii}{(ii)}\(\Rightarrow\)\SGARefLink{sga3:xvii:corollary:3:8:item:i}{(i)}, suppose that \(G\) is a unipotent
algebraic group in the sense of BIBLE, and write \(G^0\) for its
identity component.  The maximal tori of \(G^0\), being made up of
unipotent elements, are trivial; thus \(G^0\) is equal to each of its
Cartan subgroups.  It follows that \(G^0\) is solvable
(BIBLE~6, Th.~6), and hence triangularizable
(BIBLE~6, Th.~1).  In short, \(G^0\) is an algebraic subgroup of a
group \(\operatorname{Trigstr}(n)\), and is therefore unipotent in
the sense of this Exposé.

The group \((G/G^0)(k)\) is a finite group consisting of unipotent
elements.  It is therefore trivial if \(p=0\), and is a finite
\(p\)-group if \(p>0\) (BIBLE~4, Prop.~4).  But then \(G/G^0\) is
unipotent in the sense of this Exposé, since it is an iterated
extension of groups isomorphic to \(\mathbb Z/p\mathbb Z\).  This
proves that \(G\) is unipotent.

\SGAProof{Completion of the proof of \SGARefLink{sga3:xvii:theorem:3-5}{3.5}}
It remains to prove \SGARefLink{sga3:xvii:theorem:3:5:item:i}{(i)}\(\Rightarrow\)\SGARefLink{sga3:xvii:theorem:3:5:item:iv}{(iv)}.

\smallskip
\noindent\textit{\SGARefLink{sga3:xvii:appendix:A}{(a)} \(p>0\).}\quad
Consider the increasing sequence of algebraic subgroups of \(G\)
\[
  \{e\}\subset F(G)\subset F^2(G)\subset\cdots\subset
  F^n(G)\subset G^0\subset G.
\]
This gives a characteristic composition series of \(G\)
(Appendix~II, 1), and for \(n\) sufficiently large,
\(G/F^n(G)\) is smooth (Appendix~II, \SGARefLink{sga3:xvii:lemma:3-1}{3.1}).  Thus its successive
quotients are, in order:
\begin{enumerate}
\item[(1)] radicial groups of height \(1\);
\item[(2)] a smooth connected group;
\item[(3)] an étale group.
\end{enumerate}
To prove \SGARefLink{sga3:xvii:lemma:3-1:item:i}{(i)}\(\Rightarrow\)(iv), it is therefore enough to prove the
following lemma.

\SGAStatement{\SGAInlineRefTarget{sga3:xvii:lemma:3-9}Lemma 3.9}{Let \(G\) be a unipotent algebraic group
defined over a field \(k\) of characteristic \(p>0\).  Then \(G\)
has a characteristic composition series, defined over \(k\), whose
successive quotients are isomorphic to:
\begin{enumerate}
\item[(i)] \SGARefTarget{sga3:xvii:lemma:3:9:item:i}
  \((\boldsymbol\alpha_p)^r\), if \(G\) is radicial;
\item[(ii)] \SGARefTarget{sga3:xvii:lemma:3:9:item:ii}
  a twisted form of \((\mathbb G_a)^r\), if \(G\) is smooth and
  connected;
\item[(iii)] \SGARefTarget{sga3:xvii:lemma:3:9:item:iii}
  a twisted form of \((\mathbb Z/p\mathbb Z)^r\), if \(G\) is
  étale.
\end{enumerate}}
\SGARefTarget{sga3:xvii:lemma:3:9}
